% =========================================================================
% Design of a Key-Generation Cryptographic Engine for Integrated-Circuit
% Implementation: a Transistor-Level Characterized Entropy Source with
% FPGA Verification against an ESP32
%
% Master's Thesis -- Master Universitario en Microelectronica:
% Diseno y Aplicaciones de Sistemas Micro/Nanometricos
% Universidad de Sevilla / IMSE-CNM-CSIC -- Academic year 2025-2026
%
% Author: Mario Garcia Jimenez
% Compile on Overleaf: pdfLaTeX + BibTeX (main file: tfm_main.tex)
% =========================================================================
\documentclass[12pt, a4paper, oneside]{report}

\usepackage[utf8]{inputenc}
\usepackage[T1]{fontenc}
\usepackage[english]{babel}
\usepackage[margin=2.5cm]{geometry}
\usepackage{graphicx}
\usepackage{amsmath, amssymb}
\usepackage{booktabs}
\usepackage{multirow}
\usepackage{longtable}
\usepackage{array}
\usepackage{listings}
\usepackage{xcolor}
\usepackage{tikz}
\usetikzlibrary{arrows.meta, positioning, shapes.geometric, calc, chains}
\usepackage{pgfplots}
\pgfplotsset{compat=1.17}
\usepackage[hidelinks]{hyperref}
\usepackage[nottoc]{tocbibind}
\usepackage{acronym}
\usepackage{caption}
\usepackage{subcaption}
\usepackage{setspace}
\onehalfspacing

\graphicspath{{figures/}}

% ---- listings style ------------------------------------------------------
\definecolor{codegray}{rgb}{0.5,0.5,0.5}
\definecolor{codegreen}{rgb}{0,0.5,0.2}
\definecolor{backcolour}{rgb}{0.97,0.97,0.97}
\lstdefinestyle{code}{
  backgroundcolor=\color{backcolour},
  commentstyle=\color{codegreen}\itshape,
  keywordstyle=\color{blue},
  numberstyle=\tiny\color{codegray},
  stringstyle=\color{purple},
  basicstyle=\ttfamily\footnotesize,
  breaklines=true, numbers=left, numbersep=6pt,
  frame=single, framerule=0.2pt, tabsize=2, captionpos=b,
  literate={—}{{--}}1 {–}{{-}}1 {±}{{$\pm$}}1 {°}{{$^\circ$}}1
           {·}{{$\cdot$}}1 {→}{{$\rightarrow$}}1 {≈}{{$\approx$}}1
           {σ}{{$\sigma$}}1 {τ}{{$\tau$}}1 {µ}{{$\mu$}}1 {í}{{\'i}}1
           {á}{{\'a}}1 {é}{{\'e}}1 {ó}{{\'o}}1 {ú}{{\'u}}1 {ñ}{{\~n}}1,
}
\lstset{style=code}
\lstdefinelanguage{VHDL2008}[]{VHDL}{morekeywords={process,all,generic,port}}

% ---- commands ------------------------------------------------------------
\newcommand{\thetitle}{Design of a Key-Generation Cryptographic Engine for
Integrated-Circuit Implementation: a Transistor-Level Characterized Entropy
Source with FPGA Verification against an ESP32}
% Measured / simulated / estimated markers used throughout the text.
\newcommand{\meas}{\textsc{[measured]}}
\newcommand{\simu}{\textsc{[simulated]}}
\newcommand{\est}{\textsc{[estimated]}}

\begin{document}

% =========================================================================
\begin{titlepage}
\centering
{\scshape\large Universidad de Sevilla\\
Instituto de Microelectr\'onica de Sevilla (IMSE-CNM-CSIC)\par}
\vspace{0.8cm}
{\scshape M\'aster Universitario en Microelectr\'onica:\\
Dise\~no y Aplicaciones de Sistemas Micro/Nanom\'etricos\par}
\vspace{2.0cm}
{\LARGE\bfseries \thetitle\par}
\vspace{2.0cm}
{\Large Master's Thesis\par}
\vspace{1.4cm}
{\large\bfseries Author:\\ Mario Garc\'ia Jim\'enez\par}
\vspace{0.7cm}
{\large\bfseries Supervisor:\\ {[To be assigned]}\par}
\vfill
{\large Academic year 2025--2026\par}
\end{titlepage}

% =========================================================================
\chapter*{Abstract}
\addcontentsline{toc}{chapter}{Abstract}

% [DRAFT -- to be finalized with measured results]
Cryptographic keys are only as good as the entropy behind them, yet most
embedded key generators are evaluated as black boxes: statistical test
suites are run on their output and a pass is taken as evidence of
randomness. The author's previous work on the hardware random number
generator of the ESP32 showed why that is not enough --- a source with
low min-entropy can pass NIST SP~800-22 --- and NIST itself has since
rejected the suite for assessing cryptographic generators. This thesis
takes the opposite route: a white-box key-generation engine whose entropy
is justified from the physics of the device. A ring-oscillator entropy
source is designed at transistor level and its jitter is extracted from
simulation and turned, through a validated stochastic model, into a lower
bound on min-entropy per bit and into the design parameters of the
sampler. Around it, a fully portable digital core --- SP~800-90B health
tests, an AES-based conditioning and \acs{CTRDRBG} with derivation
function, AES-CMAC authentication and an \acs{I2C} register interface ---
is written once and verified three ways: against the official test
vectors, in simulation on a Basys~3 FPGA acting as the verification
vehicle, and against an ESP32 that recomputes every result with an
independent implementation (mbedtls). The same VHDL is finally
synthesized on an open 130~nm standard-cell library to give the area in
gate equivalents and the path to a dedicated integrated circuit. All
code, models, vectors and data are released.

\paragraph{Keywords} true random number generator, ring oscillator,
jitter, min-entropy, stochastic model, AIS~20/31, SP~800-90B, CTR\_DRBG,
AES-CMAC, secure element, FPGA, ASIC, ESP32.

\chapter*{Resumen}
\addcontentsline{toc}{chapter}{Resumen}

% [BORRADOR -- se cierra con los resultados medidos]
Las claves criptogr\'aficas valen lo que vale la entrop\'ia que hay
detr\'as, y sin embargo la mayor\'ia de los generadores embebidos se
eval\'uan como cajas negras: se pasan bater\'ias de tests estad\'isticos
sobre su salida y un aprobado se toma como prueba de aleatoriedad. El
trabajo previo del autor sobre el generador hardware del ESP32 mostr\'o
por qu\'e no basta --- una fuente con poca min-entrop\'ia puede pasar
NIST SP~800-22 --- y el propio NIST ha rechazado despu\'es esa bater\'ia
para evaluar generadores criptogr\'aficos. Este TFM toma el camino
contrario: un motor de generaci\'on de claves de caja blanca cuya
entrop\'ia se justifica desde la f\'isica del dispositivo. Se dise\~na a
nivel de transistor una fuente de entrop\'ia por oscilador de anillo, se
extrae su jitter por simulaci\'on y se convierte, mediante un modelo
estoc\'astico validado, en una cota inferior de min-entrop\'ia por bit y
en los par\'ametros de dise\~no del muestreador. A su alrededor, un
n\'ucleo digital totalmente portable --- tests de salud SP~800-90B,
acondicionado y \acs{CTRDRBG} con funci\'on de derivaci\'on basados en
AES, autenticaci\'on AES-CMAC e interfaz de registros \acs{I2C} --- se
escribe una sola vez y se verifica de tres maneras: contra los vectores
oficiales, en simulaci\'on sobre una FPGA Basys~3 que hace de veh\'iculo
de verificaci\'on, y frente a un ESP32 que recalcula cada resultado con
una implementaci\'on independiente (mbedtls). El mismo VHDL se sintetiza
por \'ultimo sobre una biblioteca abierta de celdas est\'andar de 130~nm
para dar el \'area en puertas equivalentes y el camino hacia un circuito
integrado dedicado. Todo el c\'odigo, modelos, vectores y datos se
publican.

\chapter*{Acknowledgements}
\addcontentsline{toc}{chapter}{Acknowledgements}
To my family, my colleagues at Airzone, and the faculty of the M\'aster en
Microelectr\'onica at IMSE-CNM, for their support during this work.

\tableofcontents
\listoffigures
\listoftables

% =========================================================================
\chapter*{List of Acronyms}
\addcontentsline{toc}{chapter}{List of Acronyms}
\begin{acronym}[CTR\_DRBG]
\acro{AES}{Advanced Encryption Standard}
\acro{AIS}{Anwendungshinweise und Interpretationen zum Schema (BSI)}
\acro{APT}{Adaptive Proportion Test}
\acro{ASIC}{Application-Specific Integrated Circuit}
\acro{BSI}{Bundesamt f\"ur Sicherheit in der Informationstechnik}
\acro{CAVP}{Cryptographic Algorithm Validation Program}
\acro{CDC}{Clock Domain Crossing}
\acro{CMAC}{Cipher-based Message Authentication Code}
\acro{CTRDRBG}[CTR\_DRBG]{Counter-mode Deterministic Random Bit Generator}
\acro{DRBG}{Deterministic Random Bit Generator}
\acro{ERO}{Elementary Ring Oscillator}
\acro{FPGA}{Field-Programmable Gate Array}
\acro{GE}{Gate Equivalent}
\acro{I2C}{Inter-Integrated Circuit}
\acro{KAT}{Known-Answer Test}
\acro{LUT}{Look-Up Table}
\acro{MEMS}{Micro-Electro-Mechanical System}
\acro{MPW}{Multi-Project Wafer}
\acro{NIST}{National Institute of Standards and Technology}
\acro{PDK}{Process Design Kit}
\acro{PLL}{Phase-Locked Loop}
\acro{PUF}{Physically Unclonable Function}
\acro{RCT}{Repetition Count Test}
\acro{RO}{Ring Oscillator}
\acro{RTL}{Register-Transfer Level}
\acro{SPI}{Serial Peripheral Interface}
\acro{TRNG}{True Random Number Generator}
\acro{XADC}{Xilinx Analog-to-Digital Converter}
\end{acronym}

% =========================================================================
\chapter{Introduction}
\label{ch:ch1_intro}

% [DRAFT v1] Source material: docs/plan.md, docs/arquitectura.md, docs/via_asic.md.

\section{Motivation: keys are only as good as their entropy}
\label{sec:motivation}

Every secure connection an embedded device opens, every firmware image it
authenticates and every credential it stores rests on a key, and every
key rests on a source of randomness. If that source is weaker than
assumed, no amount of cryptographic sophistication downstream recovers
what was lost at the start: a 128-bit key drawn from a generator with
effectively 40 bits of entropy is a 40-bit key. This is not a theoretical
concern. Frequency injection on the supply rail of a commercial secure
microcontroller collapsed its key space from $2^{64}$ to about
3300~\cite{markettos_2009_frequency}; a certified smart-card generator was
found biased despite passing the standard statistical
suites~\cite{hurleysmith_2018_certifiably}; and the hardware generator of
one of the most widely deployed IoT microcontrollers, the ESP32, silently
returns pseudo-random values when its radio is off while still passing
every black-box test~\cite{blancoromero_2026_entropy}.

The common thread is that randomness has been judged by looking at the
output. Statistical batteries such as NIST SP~800-22 check whether a bit
stream \emph{looks} uniform and independent; they cannot tell a physical
noise source from a deterministic generator seeded with a constant, and
NIST said so formally in 2022 when it announced the revision of that
publication, ``rejecting its use for assessing cryptographic random
number generators''~\cite{nist_2022_decision}. The current standards
--- AIS~20/31 in Europe~\cite{peter_2024_ais31} and SP~800-90B in the
United States~\cite{turan_2018_sp80090b} --- ask for something different:
a stochastic model of the noise source, an estimate of its min-entropy
derived from that model, and health tests that watch the source itself
rather than its output. That is a microelectronics problem before it is a
cryptography problem, and it is the problem this thesis addresses.

\section{From the black-box ESP32 study to a white-box engine}
\label{sec:background}

This work continues a previous project by the author that characterized
the hardware random number generator of the ESP32 as a black box: the
same peripheral was exercised under four operating conditions, its output
was run through a self-validated implementation of the fifteen tests of
SP~800-22, and a key-harvesting scheme combining several on-chip sources
was proposed. Its central finding was methodological: the tests passed
under conditions in which the entropy could not have been what the tests
implied, which is exactly the blind spot described above. Two of its
numbers also had to be revisited with hindsight --- a temperature
correlation was reported from a sensor whose resolution was a single
step over the observed range, and a stability figure was inferred where
it could have been measured --- and both lessons are turned into rules in
Section~\ref{sec:conventions}.

The natural next step was to build the generator instead of measuring
someone else's: an engine whose entropy source is designed, whose jitter
is measured inside the device, whose min-entropy follows from a model,
and whose keys are produced and used by primitives verified against the
standards. The ESP32 changes role from subject to judge. It has its own
AES and SHA accelerators and a mature cryptographic library, so it can
recompute everything the engine produces with an implementation that
shares nothing with it, over a real bus.

The engine is designed for silicon. Secure elements --- small integrated
circuits that generate, store and use keys on behalf of a host
microcontroller over I2C --- are a mature product category; the
\emph{transparency} of their entropy is not. Of three widely used parts,
one ships with a NIST validation of its entropy source and two do not
mention SP~800-90 anywhere in their documentation~\cite{microchip_atecc608b_ds}.
The FPGA in this work is therefore the verification vehicle, not the
target: it lets millions of test vectors run against the design in
minutes, which no transistor-level simulator can, before the same
description is synthesized on standard cells.

\section{Objectives}
\label{sec:objectives}

\begin{enumerate}
\item[O1] Design a ring-oscillator entropy source at transistor level,
with enable and sampler, and simulate it with noise, stating the
numerical resolution of the bench before any figure is read from it.
\item[O2] Extract the jitter from the simulated waveforms with an
estimator that separates quantization, thermal and flicker components,
validated beforehand on synthetic data of known jitter.
\item[O3] Turn the jitter into a min-entropy bound through a stochastic
model checked against an exact computation, and derive the sampler
divider that meets the AIS~20/31 target with margin.
\item[O4] Write the digital engine --- health tests, AES-based
conditioning and CTR\_DRBG, AES-CMAC and an I2C register interface --- as
technology-independent RTL verified against official vectors.
\item[O5] Verify the engine functionally on an FPGA with the ESP32 as an
independent oracle: known-answer tests, cross interoperability and
challenge-response.
\item[O6] Run an entropy campaign on real silicon (the FPGA's rings) with
the SP~800-90B estimators and compare with the ESP32 generator under the
same methodology.
\item[O7] Synthesize the engine on an open standard-cell library and
report its area in gate equivalents with the growth factors to core area.
\item[O8] Compare the design with the ESP32 generator, with commercial
secure elements and with published circuits, with stated uncertainties.
\end{enumerate}

The minimum defensible result is O1--O3 together with O5: an entropy
source designed from the transistor, with its jitter extracted and turned
into a bound, and a digital engine verified against an independent
implementation. A negative result that is well measured closes the work
as well as a positive one: if the accumulation exponent on hardware is
two rather than one, flicker dominates and the divider cannot grow
indefinitely, and that is a finding.

\section{Scope and exclusions}
\label{sec:scope}

Four things are deliberately out of scope, each for a stated reason.
\emph{Fabrication}: an academic shuttle is affordable and the engine fits
in the smallest block offered, but the schedule from tape-out to packaged
and characterized samples exceeds a master's thesis; the cost and the path
are documented, not executed. \emph{Asymmetric cryptography}: an
elliptic-curve multiplier is a project in itself and would consume the
whole area budget; the engine provides symmetric keys and authentication.
\emph{Side-channel resistance}: the threat model is a logical, remote
adversary; timing leaks are avoided by construction, but no masking or
trace-based evaluation is attempted, and the institute's own laboratory
for such work is noted as the natural continuation. \emph{Formal
certification}: the criteria of AIS~20/31 and SP~800-90B and the official
vectors are used throughout, and no validation is claimed.

\section{Methodology and reproducibility conventions}
\label{sec:conventions}

Four conventions apply to every figure in this document. First, each
number is labelled \meas{}, \simu{} or \est{}, and carries its
uncertainty or the source it comes from. Second, \textbf{the resolution of
the instrument is declared before any trend is interpreted}, whether the
instrument is a sensor, an estimator or a simulator --- the numerical
noise floor of the SPICE bench and the quantization of a temperature
sensor are treated alike. Third, quantities are measured rather than
inferred whenever a direct measurement exists. Fourth, novelty is stated
as ``not found in the literature consulted'', never as ``does not exist'';
one such claim in an early draft of this work had to be corrected when a
2024 paper measuring exactly the quantity in question was found, and the
correction is left visible in the working notes.

Everything is released: the VHDL and its testbenches, the Python models
and analysis, the firmware, the synthesis flow, the official test vectors
used and the raw data of every campaign, with the commands that reproduce
each table.

\section{Structure of the document}
\label{sec:structure}

Chapter~\ref{ch:ch2_soa} reviews ring-oscillator generators, their
stochastic models, the standards and the commercial landscape.
Chapter~\ref{ch:ch3_entropy} builds the chain from transistor to
min-entropy and validates every link. Chapter~\ref{ch:ch4_engine}
describes the digital engine and Chapter~\ref{ch:ch5_verification} its
verification. Chapter~\ref{ch:ch6_asic} synthesizes the engine on standard
cells and states the path and cost of an integrated circuit.
Chapter~\ref{ch:ch7_results} reports the hardware campaign and the
comparisons, and Chapter~\ref{ch:ch8_concl} concludes.

\chapter{State of the Art}
\label{ch:ch2_soa}

% [DRAFT v1] Source material: docs/estado_del_arte_trng.md, estado_del_arte_cripto.md,
% plataforma_y_flujo.md, via_asic.md. Every numeric claim below carries a
% reference that was checked against its source when the review was compiled;
% items the review could not open are marked as such in the working notes.

The review is organized along the chain the thesis builds: how ring
oscillators are turned into random bits, how their entropy is modeled and
bounded, how their jitter is measured from inside the device, what the two
governing standards demand, which symmetric primitives fit a constrained
engine, what the commercial parts offer, and what changes when the design
leaves the FPGA for silicon. The last section states the gap this work
addresses, in the guarded form the working rules require.

\section{Ring-oscillator TRNG architectures}
\label{sec:soa_archs}

A ring oscillator is an odd chain of inverters closed on itself; its
period fluctuates because each transition is delayed by a random amount,
and that fluctuation --- jitter --- is the noise source of an entire family
of generators. What distinguishes the architectures is how the jitter is
accumulated and sampled~\cite{petura_2016_survey, lubicz_2024_recommendations}.

The \textbf{elementary ring oscillator} (ERO) samples one ring with a
flip-flop clocked by a second ring divided by $K_D$; the divider sets the
accumulation time. It has the most mature stochastic
model~\cite{baudet_2011_security, killmann_2008_design, ma_2014_entropy} and
an embedded jitter measurement~\cite{fischer_2014_embedded}, at the price
of a very low output rate ($K_D$ of $10^4$ to $10^5$: 80\,000 in Spartan-6
and 135\,000 in Cyclone~V in Petura's implementation, about 430\,000 in
Fischer and Lubicz's design for 0.997~bit of entropy). The
\textbf{multi-ring} generator of Sunar, Martin and
Stinson~\cite{sunar_2007_provably} XORs many rings sampled by a reference
clock; Wold and Tan~\cite{wold_2009_analysis} added a flip-flop per ring
before the XOR and reported 100~Mbit/s in under 100 logic elements. Its
weakness was exposed by Bochard \emph{et al.}~\cite{bochard_2010_true}:
part of the apparent randomness is \emph{pseudo}-randomness from the
interaction of rings at different frequencies, and the independence the
proof assumes does not hold on an FPGA. \textbf{Coherent sampling} (COSO)
uses two rings of nearly equal period and counts the beat~\cite{kohlbrenner_2004_embedded,
valtchanov_2009_enhanced}; its model gives min-entropy directly from the
distribution of the count~\cite{peetermans_2021_configurable}, and
dynamic calibration of reconfigurable rings makes it portable across
devices and placements (3.30~Mbit/s in Spartan-6, 4.65~Mbit/s in
Spartan-7)~\cite{peetermans_2019_portable, peetermans_2021_configurable}.
The \textbf{transition-effect ring} (TERO) counts how many times a
bistable structure oscillates before settling~\cite{varchola_2010_new},
with a full physical-to-stochastic model~\cite{bernard_2019_physical};
the \textbf{self-timed ring} (STR) circulates several events in an
asynchronous ring to obtain phases spaced by the order of the
jitter~\cite{cherkaoui_2013_very}, reaching 154~Mbit/s; and the
\textbf{ES-TRNG} of Yang \emph{et al.} uses the carry chain as a
time-to-digital converter to encode the edge position with 10~LUTs and
5~flip-flops~\cite{yang_2018_estrng, rozic_2015_highly}.

Table~\ref{tab:petura} reproduces the comparison Petura \emph{et al.}
made of the AIS-compliant cores on one device~\cite{petura_2016_survey}.
The entropy column is a statistical estimate (the AIS~31 T8 test), not a
model bound; the feasibility column is the authors' qualitative note,
where COSO and TERO score worst because they need manual adjustment per
device. No comparable table exists for Artix-7; the open OpenTRNG project
implements ERO, multi-ring and COSO on an Arty~A7 with the same XC7A35T
die as the Basys~3 but publishes neither entropy nor throughput
figures~\cite{opentrng_2025_docs}.

\begin{table}[htbp]
\centering
\caption{AIS-compatible TRNG cores on Xilinx Spartan-6 (45~nm), after
Petura \emph{et al.}~\cite{petura_2016_survey}, Table~II. Entropy is the
T8 estimate; ``feasibility'' is the authors' rank (1 best).}
\label{tab:petura}
\begin{tabular}{@{}lrrrrl@{}}
\toprule
Core & LUT/reg. & Power (mW) & Rate (Mbit/s) & Entropy/bit & Feasib. \\
\midrule
ERO  & 46/19   & 2.16  & 0.0042 & 0.999 & 5 \\
COSO & 18/3    & 1.22  & 0.54   & 0.999 & 1 \\
MURO & 521/131 & 54.72 & 2.57   & 0.999 & 4 \\
PLL  & 34/14   & 10.6  & 0.44   & 0.981 & 3 \\
TERO & 39/12   & 3.31  & 0.625  & 0.999 & 1 \\
STR  & 346/256 & 65.9  & 154    & 0.998 & 2 \\
\bottomrule
\end{tabular}
\end{table}

\section{Stochastic models and entropy bounds}
\label{sec:soa_models}

The requirement that a physical generator come with a \emph{stochastic
model} --- a mathematical description of the noise source from which the
entropy of the output can be derived --- was introduced by Killmann and
Schindler~\cite{killmann_2008_design, schindler_2002_evaluation} and is
mandatory in AIS~20/31. For ring oscillators the reference model is
Baudet \emph{et al.}~\cite{baudet_2011_security}: the phase of the sampled
ring is a Wiener process with drift $\mu$ and volatility $\sigma^2$, the
sampled bit is the half-period the phase falls in, and everything depends
on the quality factor $Q=\sigma^2\Delta t$ (accumulated phase variance
between samples) and the frequency ratio $\nu=\mu\Delta t$. The model
gives the bit probability conditioned on the previous phase, the
probability of any output vector, and a lower bound on the Shannon
entropy rate, $H \ge 1 - \frac{4}{\pi^2\ln 2}e^{-4\pi^2 Q}$, whose exact
status is examined in Chapter~\ref{ch:ch3_entropy}. Petura \emph{et al.}
rewrite it in design quantities, $H = 1 - \frac{4}{\pi^2\ln2}\exp(-\pi^2
\sigma_{th}^2 K T_2/T_1^3)$, and Fischer and Lubicz invert it to obtain
the divider $K_D$ from a measured jitter~\cite{fischer_2014_embedded}. Ma
\emph{et al.} give an alternative model used by the Leuven
group~\cite{ma_2014_entropy}; Haddad \emph{et al.} question the mutual
independence of jitter realizations on which these cancellations
rest~\cite{haddad_2014_assumption}; and Saarinen warns that Baudet's bound
``is never lower than 0.415 even when $Q$ approaches zero'' and is safe
only under additional assumptions, proposing a bound valid over the whole
range~\cite{saarinen_2021_entropy}.

Two refinements matter for a design. First, not all jitter is entropy.
Fischer and Lubicz separate the accumulated phase variance into a
random-walk term, linear in time, produced by thermal (white) noise, and
flicker ($1/f^\beta$) terms whose variance grows quadratically and
dominate at long accumulation times~\cite{fischer_2014_embedded};
deterministic, global jitter from the supply and substrate behaves like
the latter and can be known or induced by an
adversary~\cite{fischer_2008_enhancing}. Only the linear term may be
credited, which is why a divider cannot be made arbitrarily large to
compensate a weak source. Second, every embedded measurement is
\emph{differential}: two rings sharing a supply see the same global
jitter, which cancels in their relative phase, leaving the sum of their
local jitters. Lubicz and Sk\'orski formalize this ``jitter transfer
principle'' and show that with two rings one obtains only the sum, while
three rings and three pairwise measurements recover the individual
variances; their experiment also shows the assumption of variance linear
in period failing by a factor of two in one case~\cite{lubicz_2024_jitter}.

\section{Embedded jitter measurement}
\label{sec:soa_jitter}

Oscilloscopes cannot resolve the picoseconds of jitter of a ring inside a
chip without an external divider~\cite{benea_2024_flicker}, so the useful
methods measure from within. The \textbf{counter method} counts periods of
one ring in a fixed window of the other; the variance of the count grows
with the accumulated relative jitter~\cite{killmann_2008_design,
lubicz_2015_towards}, with a quantization floor of one count and
contamination by flicker at long windows. The method of \textbf{Fischer
and Lubicz}~\cite{fischer_2014_embedded} needs no divider: in $K\approx10^4$
blocks of $N\approx100$ bits sampled at $K_D=1$, the fraction of pairs
$(b_j,b_{j+M})$ that differ tracks the phase accumulated over $M$ periods,
and its variance across blocks, plotted against $M$ between 200 and 1600,
gives the jitter per period with an error under 5\,\% in simulation; on a
Cyclone~III they measured $\sigma=5.01$~ps per 7.81~ns period, and showed
that cooling the chip reduces $\sigma$ --- which makes the same measurement
an online test. \textbf{Coherent sampling} yields the accumulated jitter
from the variance of the beat count with no extra
hardware~\cite{peetermans_2021_configurable}. Most recently Benea
\emph{et al.} characterize the jitter of rings with the \textbf{Allan
variance}, fitting $\sigma_t^2(t)=a_0+a_1 t+a_2 t^2$ to separate
quantization, thermal and flicker terms, and do so both on an Artix-7
--- the FPGA family of this work --- and on a 28~nm FD-SOI circuit, with
data and scripts released~\cite{benea_2024_flicker}. Published per-period
jitter figures cluster at 2--4~ps for periods of 4--8~ns in Spartan-6 and
Cyclone~V~\cite{petura_2016_survey}, 5~ps in Cyclone~III measured
embedded~\cite{fischer_2014_embedded}, and 2--2.5~ps in Cyclone~III and
Virtex-5 for self-timed rings~\cite{cherkaoui_2013_very}; the Artix-7
coefficients of Benea \emph{et al.} are published in counter units, not
seconds, so an absolute figure for the XC7A35T remains to be measured.

\section{Standards: AIS 20/31 v3.0 and NIST SP 800-90B}
\label{sec:soa_standards}

Two documents govern how a physical generator is assessed. The German
\textbf{AIS~20/31}, in its version~3.0 of September 2024~\cite{peter_2024_ais31},
replaces the 2011 requirement of 0.997~bit of Shannon entropy per
internal bit~\cite{killmann_2011_ais31} with a choice in class PTG.2:
Shannon entropy $\ge0.9998$ and/or min-entropy $\ge0.98$, together with
$P(1)\in(0.493, 0.507)$; the BSI's transition policy allows new
evaluations under the old version only until March
2026~\cite{bsi_2025_transition}. Its requirements are structural as much
as numerical: a verifiable stochastic model is mandatory; the online test
must be derived from the model, failing with high probability when the
source parameter leaves the admissible set; a total-failure test must
detect entropy collapsing to zero, and the repetition-count test of
SP~800-90B is explicitly accepted for it; and the black-box suite
$T_{\mathrm{irn}}$ now includes two of NIST's predictors (MultiMMC and
LZ78Y), a deliberate harmonization~\cite{peter_2024_ais31, schindler_2024_ecw}.

The American \textbf{SP~800-90B}~\cite{turan_2018_sp80090b} estimates
\emph{min-entropy} --- the adversarial quantity, $-\log_2$ of the most
probable value --- from at least one million raw samples and a
1000$\times$1000 restart matrix, using the most-common-value estimator
if the source is shown to be i.i.d.\ and otherwise the minimum of ten
estimators, four of them predictors that try to guess the next sample.
It approves two continuous health tests and requires no others: the
Repetition Count Test with cutoff $C=1+\lceil-\log_2\alpha/H\rceil$ and the
Adaptive Proportion Test over a window of 1024 samples for binary sources
(512 for non-binary), $C=1+\mathrm{CRITBINOM}(W,2^{-H},1-\alpha)$, with
$\alpha=2^{-20}$ recommended. It lists the conditioning components whose
output may claim full entropy --- HMAC, CMAC, CBC-MAC, approved hashes and
the two derivation functions of SP~800-90A --- and gives the formula for
the entropy of a conditioned output as a function of input length, input
entropy and internal width. The reference implementation of the
estimators is public~\cite{nist_ea_github}.

The two standards agree on what the older statistical suites cannot do.
SP~800-22~\cite{bassham_2010_sp80022} tests uniformity and independence of
a bit stream; a deterministic generator with zero entropy passes it by
construction, and so does a weak physical source behind a cryptographic
post-processor~\cite{saarinen_2021_entropy}. NIST's 2022 decision to
revise the publication states the point in so many words, ``rejecting its
use for assessing cryptographic random number generators''~\cite{nist_2022_decision}.
The empirical record agrees: a certified smart-card generator showed
consistent biases while passing most tests~\cite{hurleysmith_2018_certifiably,
hurleysmith_2018_great}, and the ESP32's generator passes the same
screenings with its radio off, when Espressif's own documentation says it
returns pseudo-random values~\cite{espressif_idf_random, espressif_trm_esp32,
blancoromero_2026_entropy}. The academic literature on that generator
applies only SP~800-22 or Dieharder~\cite{castillo_2026_automated,
hidayatulloh_2024_performance}; no SP~800-90B characterization of it was
found.

\section{Symmetric cores for constrained hardware}
\label{sec:soa_cores}

The engine needs a block cipher and, if it uses one, a hash. For AES-128
(ten rounds~\cite{nist_2023_fips197}) the design space runs from
fully unrolled pipelines for throughput to 8-bit datapaths for area:
Good and Benaissa report 124 slices and two block RAMs at 2.2~Mbit/s on
a Spartan-II at one extreme and 25~Gbit/s at the
other~\cite{good_2005_fastest}; Chodowiec and Gaj's compact design uses 222
slices and three block RAMs~\cite{chodowiec_2003_compact}; and the
smallest ASIC realizations reach 2.4~kGE with an 8-bit datapath at 226
cycles per block~\cite{moradi2011pushing}, following the composite-field
S-box line of Satoh~\cite{satoh_2001_compact}, Canright~\cite{canright_2005_compact}
and Boyar--Peralta~\cite{boyar_2010_minimization}. The composite-field
constructions optimize ASIC gate counts and lose their advantage in an
FPGA LUT, where a direct table is both faster and simpler. Open
reference designs with published Artix-7 figures exist: a one-round-per-cycle
encryptor at 1104~LUTs and 294~MHz~\cite{hadipour_aesvhdl}, and a 32-bit
iterative encryptor/decryptor at 2298 slices and 97~MHz~\cite{strombergson_aes}.
For an engine whose output leaves through an I2C link, any of these is
far faster than the link, and the choice is dictated by verifiability and
area, not speed.

SHA-256~\cite{nist_2015_fips1804} in its iterative form costs about 750
slices at 108~MHz on Artix-7~\cite{strombergson_sha256}, comparable to a
compact AES; its role in a generator would be as the vetted conditioner of
SP~800-90B. The deterministic generators of
SP~800-90A~\cite{barker_2015_sp80090ar1} --- Hash\_DRBG, HMAC\_DRBG and
CTR\_DRBG --- differ mainly in cost: CTR\_DRBG with AES-128 keeps a state
of 256~bits plus a counter and produces 128~bits per block encryption,
Hash\_DRBG needs 440-bit modular additions, and HMAC\_DRBG two hash
compressions per output block. Without a derivation function CTR\_DRBG
requires full-entropy input of exactly the seed length; with one, the
usual configuration in practice, it accepts arbitrary input. Woodage and
Shumow's analysis warns that the standard's flexibility permits choices
that worsen state leakage and recommends updating after every generate
and reseeding far below the $2^{48}$ limit~\cite{woodage_2019_analysis};
SP~800-90C, final since 2025, frames how a physical source and a DRBG
combine~\cite{barker_2025_sp80090c}. For authentication, AES-CMAC
(SP~800-38B~\cite{nist_2016_sp80038b}, RFC~4493~\cite{song_2006_rfc4493})
costs one block encryption per message block and is itself a vetted
conditioner. The official vectors for all of these are
public~\cite{nist_cavp_blockciphers, nist_cavp_modes, nist_cavp_drbg,
nist_acvp_server}.

\section{Commercial secure elements}
\label{sec:soa_se}

The product category this engine belongs to exists. A secure element is
a small integrated circuit that generates and stores keys and performs
symmetric and asymmetric operations for a host microcontroller, almost
always over I2C. Four parts dominate the low-cost segment: Microchip's
ATECC608B, NXP's SE050, Infineon's OPTIGA Trust~M and ST's STSAFE-A110~\cite{microchip_atecc608b_ds,
nxp_se050_ds, infineon_optiga_trustm_ds, st_stsafe_a110_ds}. The
ATECC608B, at 0.90~\$ in single units~\cite{digikey_prices_2026}, offers
AES-128, SHA-256/HMAC, ECDSA/ECDH on P-256, sixteen key slots and an I2C
interface, and its entropy source holds a NIST validation under
SP~800-90B, obtained precisely with ring oscillators~\cite{nist_esv_e46}.
Two of its competitors make no reference to SP~800-90 in their datasheets;
the third states that its generator is ``designed to meet'' the standard.
The lesson for this work is where an academic design can and cannot
compete: not on price, where a certified part costs less than a euro, but
on the transparency of the entropy, which is exactly what a black-box part
cannot offer and what the standards now demand.

\section{From FPGA to ASIC: what changes in the entropy source}
\label{sec:soa_asic}

Digital logic ports from an FPGA to standard cells with a known penalty
--- Kuon and Rose measured 40$\times$ in area, 3.2$\times$ in delay and
12$\times$ in dynamic power for logic-only designs in 90~nm, and 27--49$\times$
in area for AES specifically~\cite{kuon2006}; the ring oscillator does not
port at all, because in silicon it is built from inverter cells whose
noise, variability and coupling differ from a LUT chain. Three published
findings define what a silicon implementation must add. Benea
\emph{et al.} show that flicker noise, whose weight grows in newer nodes,
changes the character of the jitter and, against expectation, that
treating it as a legitimate source can raise throughput fourfold at equal
entropy~\cite{benea_2024_flicker}. Yang, Blaauw and Sylvester show that
process variation --- larger in 40~nm than in 180~nm --- must be absorbed
by a configurable ring with a tuning loop for the generator to survive
$-40$ to 120~$^\circ$C and 0.6 to 0.9~V~\cite{yang_2016_jssc}. And two
attacks show why a security device cannot ignore its environment: a
frequency injected on the supply locks the rings and removed the jitter
of a commercial secure microcontroller, reducing its key space from
$2^{64}$ to 3300~\cite{markettos_2009_frequency}, and a few hundred
microwatts of radiated field biased a fifty-ring generator until it
failed every test~\cite{bayon_2012_contactless}; the published defence
combines supply filtering with the same calibration loop~\cite{yang_2016_jssc},
and device aging under bias-temperature instability and hot-carrier
injection degrades poorly-timed generators over years while leaving
well-timed ones intact.

Access to silicon is more open than it appears. The University of Sevilla
and IMSE-CNM are members of Europractice~\cite{europractice_memberlist},
whose academic price list is public: a mini@sic block costs 3\,430~\texteuro{}
in UMC~180~nm and 3\,691~\texteuro{} in TSMC~65~nm, and IHP's SG13G2 is
offered at 5\,110~\texteuro{}/mm$^2$ with a free run for open-source designs
under 2~mm$^2$~\cite{europractice_prices2026, europractice_minisic,
ihp_openpdk}. The institute's hardware-security area has already taken
ring-oscillator primitives from Xilinx 7-series FPGAs, evaluated with
SP~800-22 and SP~800-90B, to characterized 65~nm silicon~\cite{imse_area_hwsec,
martinez2022efficientropuf, ortega2024vlsi65nm, ortega2026ropuf65nm}, and
has fabricated cryptographic circuits before~\cite{mora2020trivium}. The
open standard-cell libraries that make a synthesis chapter possible ---
IHP SG13G2, SkyWater SKY130, GlobalFoundries GF180MCU and the predictive
Nangate 45~nm --- come with liberty, LEF and GDS views, and the literature
compares such designs in gate equivalents normalized to the minimum NAND2
of the library~\cite{ihp_openpdk_repo, skywater_pdk_cells, gf180mcu_pdk,
nangate45, kaps2019lwc}, with the caution that logic-synthesis area is a
lower bound that grows by a factor of 1.7 to 4.6 to routed core area.

\section{Gap addressed by this work}
\label{sec:gap}

No absolute novelty is claimed. The ERO and its model, the embedded
jitter measurement, the survey of AIS-compliant cores, the calibrated
COSO on 7-series devices and an open implementation on the very die of
the Basys~3 all exist, and the institute's own group has crossed from
7-series FPGAs to 65~nm silicon with ring-oscillator primitives. What was
\emph{not found} in the literature consulted, and what this work sets out
to provide in a modest but verifiable way, is the following.

\begin{enumerate}
\item A per-period jitter figure for the Artix-7 XC7A35T in absolute
units, obtained inside the device with a differential method, with the
thermal and flicker components separated and with stated uncertainty ---
the published Artix-7 coefficients being in counter units.
\item The complete chain \emph{measured jitter $\to$ quality factor $\to$
min-entropy bound $\to$ divider $\to$ health-test cutoffs}, applied to one
platform and confronted with the ten non-i.i.d.\ estimators of SP~800-90B
and with the PTG.2 requirements of AIS~20/31 \emph{version~3.0}, which
most of the earlier literature could not use.
\item An SP~800-90B characterization of the ESP32 generator in its
operating states, and a controlled comparison between a white-box FPGA
source and a black-box commercial one under the same methodology.
\item The use of a commodity microcontroller with its own cryptographic
accelerators as an independent oracle for an FPGA engine, over a real
bus, with known-answer, interoperability and challenge-response tests.
\item A single-primitive key-generation engine --- AES for conditioning,
generation and authentication --- written as portable RTL, verified
against official vectors, synthesized on an open 130~nm library with the
result stated as a lower bound, and costed for an academic shuttle.
\end{enumerate}

Each item is formulated as ``not found'' rather than ``does not exist''.
One earlier claim of this kind in the working notes --- that no Artix-7
jitter measurement had been published --- had to be corrected when
Benea \emph{et al.}'s paper was found, and the correction is kept visible
there.

\chapter{Entropy Source: from Transistor to Min-Entropy}
\label{ch:ch3_entropy}

% [DRAFT v1 -- theory and simulation only; hardware figures pending]
% Source material: docs/modelo_ero_resultados.md, docs/arquitectura.md,
% analysis/ero_model.py, analysis/jitter_estimator.py, analysis/spice_ring.py

This chapter builds the chain that the rest of the thesis rests on: from
the physical jitter of a ring oscillator to a defensible figure of
min-entropy per output bit, and from that figure to the design parameters
of the sampler. Every step is validated before it is used. The
statistical estimators are validated against synthetic data of known
jitter; the analytic bounds are validated against an exact numerical
computation; and the transistor-level bench is validated by measuring its
own resolution before anything is measured with it. All figures in this
chapter are \simu{} unless stated otherwise; the hardware campaign is
reported in Chapter~\ref{ch:ch7_results}.

\section{The elementary ring oscillator}
\label{sec:ero}

The entropy source is an \ac{ERO}: a free-running ring oscillator RO1
whose output is sampled by a flip-flop clocked by a second ring
oscillator RO2 divided by an integer $K_D$~\cite{baudet_2011_security,
fischer_2014_embedded}. Between two samples the phase of RO1 accumulates
jitter for $K_D$ periods of RO2; when the accumulated jitter is
comparable to the period of RO1, the sampled bit becomes unpredictable.
The \ac{ERO} was chosen over higher-throughput alternatives (multi-ring
XOR, coherent sampling, transition-effect rings) for one reason: it is
the architecture for which a stochastic model with a provable entropy
bound exists and for which an embedded jitter measurement method has
been published and validated~\cite{petura_2016_survey}. Throughput is
not a requirement here --- the engine consumes entropy only to seed a
deterministic generator --- whereas a demonstrable bound is.

A design decision that follows from the platform is that \emph{the board
oscillator is never used as a reference}. The Basys~3 clock source is a
MEMS oscillator with a fractional \ac{PLL} whose datasheet specifies a
cycle-to-cycle jitter of up to 95~ps \est{}~\cite{digilent_2014_basys3rm},
one to two orders of magnitude above the per-period jitter of a ring
oscillator in a 28~nm process (2--4~ps \est{}~\cite{petura_2016_survey}).
Sampling RO1 with that clock would measure the phase noise of the PLL,
a global and potentially inducible component, rather than the thermal
noise of the ring. RO2 therefore samples RO1, and the 100~MHz system
clock only drives the synchronous logic.

\section{Stochastic model and the quality factor}
\label{sec:model}

Following Baudet \emph{et al.}~\cite{baudet_2011_security}, the phase
$\varphi(t)$ of RO1, measured in periods, is modeled as a Wiener process
with drift $\mu$ and volatility $\sigma^2$. Between two samples separated
by $\Delta t$ the phase advances $\nu = \mu\,\Delta t$ periods on average
and accumulates a variance
\begin{equation}
  Q = \sigma^2\,\Delta t ,
  \label{eq:Q}
\end{equation}
the \emph{quality factor}. The output bit is $s = 1$ when the fractional
phase lies in $[1/2, 1)$. Conditioned on the previous phase $\varphi_0$,
the bit is Bernoulli with
\begin{equation}
  P(s=1 \mid \varphi_0) = \sum_{k\in\mathbb{Z}}
  \left[\Phi\!\left(\tfrac{k+1-\varphi_0-\nu}{\sqrt{Q}}\right) -
        \Phi\!\left(\tfrac{k+\tfrac12-\varphi_0-\nu}{\sqrt{Q}}\right)\right],
  \label{eq:puno}
\end{equation}
where $\Phi$ is the standard normal distribution function; this is the
wrapped Gaussian evaluated exactly. Its first harmonic gives the maximum
bias $\varepsilon_{\max} \approx (2/\pi)\,e^{-2\pi^2 Q}$, from which
Baudet's Corollary~1 yields the well-known expression
\begin{equation}
  H(s \mid \varphi_0) = 1 - \frac{4}{\pi^2 \ln 2}\, e^{-4\pi^2 Q}
  + O\!\left(e^{-6\pi^2 Q}\right).
  \label{eq:baudet14}
\end{equation}

\section{Exact conditional entropy vs.\ the asymptotic bound}
\label{sec:exact}

Equation~\eqref{eq:baudet14} is routinely cited as a \emph{lower bound}
on entropy per bit. Read carefully, the inequality in the corollary is
only that the entropy rate of the sequence is at least the entropy of
one bit conditioned on the exact previous phase; the right-hand side is
an asymptotic expansion \emph{of} that conditional entropy, not a bound
on it. Expanding the binary entropy in the bias,
$h_2(\tfrac12+\varepsilon) = 1 - \tfrac{1}{\ln 2}\sum_{n\ge1}
\tfrac{(2\varepsilon)^{2n}}{2n(2n-1)}$, and averaging over a uniform
$\varphi_0$ with $\varepsilon = A\sin(2\pi\varphi)$,
$A = (2/\pi)e^{-2\pi^2 Q}$, gives
\begin{equation}
  H(s\mid\varphi_0) = 1 - \frac{4}{\pi^2\ln 2}\,e^{-4\pi^2 Q}
                        - \frac{8}{\pi^4\ln 2}\,e^{-8\pi^2 Q} - \cdots
  \label{eq:orden2}
\end{equation}
All omitted terms are negative, so \emph{any} truncation overestimates
the entropy; the next term decays as $e^{-8\pi^2 Q}$, not $e^{-6\pi^2 Q}$.
Table~\ref{tab:exact} compares the exact value computed from
\eqref{eq:puno} on a grid of 8192 phases with the first- and second-order
expansions. Adding the second-order term reduces the error of
\eqref{eq:baudet14} by a factor of 95 at $Q=0.10$ and 4970 at $Q=0.20$,
which confirms the derivation.

\begin{table}[htbp]
\centering
\caption{Conditional Shannon entropy per bit: exact wrapped-Gaussian
computation vs.\ truncated expansions \simu{}.}
\label{tab:exact}
\begin{tabular}{@{}lcccc@{}}
\toprule
$Q$ & Exact & Order 2 & Order 1 (Eq.~\ref{eq:baudet14}) & Error of order 1 \\
\midrule
0.02 & 0.7013194 & 0.7100952 & 0.7345213 & $+3.3\times10^{-2}$ \\
0.05 & 0.9163001 & 0.9164920 & 0.9187783 & $+2.5\times10^{-3}$ \\
0.10 & 0.9886728 & 0.9886733 & 0.9887174 & $+4.5\times10^{-5}$ \\
0.15 & 0.9984319 & 0.9984319 & 0.9984327 & $+8.5\times10^{-7}$ \\
0.20 & 0.9997823 & 0.9997823 & 0.9997823 & $+1.6\times10^{-8}$ \\
\bottomrule
\end{tabular}
\end{table}

In the operating region ($Q\approx0.2$) the discrepancy is $10^{-8}$ and
irrelevant; for an under-dimensioned generator ($Q\approx0.02$) the
truncated formula credits 3.3\,\% of entropy that does not exist. This
thesis therefore uses the exact computation as the conservative value and
cites \eqref{eq:baudet14} as what it is: an excellent asymptotic
approximation for $Q \gtrsim 0.1$. The same numerical machinery shows
that the min-entropy obtained from the first-harmonic bias,
$H_\infty = -\log_2(\tfrac12 + \varepsilon_{\max})$, agrees with the exact
worst-case value to better than $10^{-4}$~bit over $Q\in[0.05, 0.30]$.

\subsection{Design target}

AIS~20/31 version~3.0~\cite{peter_2024_ais31} admits two alternative
criteria for class PTG.2: Shannon entropy $\ge 0.9998$ or min-entropy
$\ge 0.98$ per internal bit. Solving both by bisection on the exact
entropies gives $Q \ge 0.2022$ and $Q \ge 0.2286$ respectively. The
min-entropy criterion is the more demanding one and is also the quantity
NIST SP~800-90B works with, so it is adopted here: \textbf{the design
target is $Q \ge 0.23$, with a set point of $Q = 0.30$} ($H_\infty =
0.9951$), the margin covering the reduction of thermal jitter at low
temperature, which is the standard adversarial test of Fischer and
Lubicz~\cite{fischer_2014_embedded}.

\section{Embedded jitter estimator and its two corrections}
\label{sec:estimator}

To choose $K_D$ one must know $\sigma^2$, and the method of Fischer and
Lubicz~\cite{fischer_2014_embedded} measures it from the bits themselves:
with $K_D = 1$, the fraction $c$ of pairs $(b_j, b_{j+M})$ that differ
within a block of $N$ pairs estimates the phase accumulated over $M$
samples, and its variance across $K$ blocks grows as $4 M Q_{\mathrm{raw}}$.
Implemented literally, with a straight-line fit, the estimator
\emph{underestimates} $Q_{\mathrm{raw}}$ by 13\,\% to 38\,\% depending on
the jitter level. Two causes were identified on synthetic data with known
$Q$ and corrected:

\begin{enumerate}
\item \textbf{Vertex effect.} $c$ is not linear in the accumulated phase;
it is a \emph{triangular} wave of period one and slope $\pm2$. When the
deterministic phase $M\nu \bmod 1$ falls near a vertex the distribution
straddles it and the variance collapses --- to 0.38 of its nominal value
in one measured case. The vertex is detected from the measured mean
itself ($\bar c \approx 0$ or $1$ is a vertex, $\bar c \approx \tfrac12$
the middle of the ramp), so distances $M$ with $\bar c$ outside
$[0.15, 0.85]$, or for which $3\sqrt{MQ}$ reaches the vertex, are
discarded iteratively.
\item \textbf{Non-constant estimation noise.} Estimating a fraction from
$N$ pairs adds a variance that scales with $\bar c(1-\bar c)$, which
changes with $M$. Fitting it as a constant intercept biases the slope; it
is fitted instead as a second regressor,
$\mathrm{Var}(c) = Q\,(4M) + \kappa\,\bar c(1-\bar c)/N$.
\end{enumerate}

After both corrections the estimator recovers the injected $Q_{\mathrm{raw}}$
within $\pm2\,\%$ over a range of $150\times$ (Table~\ref{tab:estimator}),
with a bias of $-0.6\,\%$ and a dispersion of 4.0\,\% over eight
independent seeds for $K=20\,000$ blocks. The fitted accumulation
exponent is $1.009$ for a simulator without flicker noise; measured on
hardware, the same exponent separates the thermal component ($b=1$, which
counts as entropy) from flicker and global jitter ($b=2$, which does not).

\begin{table}[htbp]
\centering
\caption{Recovery of a known $Q_{\mathrm{raw}}$ by the corrected embedded
estimator ($K=20\,000$ blocks of $N=100$ pairs) \simu{}.}
\label{tab:estimator}
\begin{tabular}{@{}cccc@{}}
\toprule
$Q_{\mathrm{raw}}$ injected & Estimated & Error & $M$ values used \\
\midrule
$2.0\times10^{-7}$ & $1.979\times10^{-7}$ & $-1.07\,\%$ & 19 \\
$5.0\times10^{-7}$ & $5.060\times10^{-7}$ & $+1.21\,\%$ & 18 \\
$2.0\times10^{-6}$ & $2.037\times10^{-6}$ & $+1.86\,\%$ & 16 \\
$8.0\times10^{-6}$ & $8.039\times10^{-6}$ & $+0.48\,\%$ & 9 \\
$3.0\times10^{-5}$ & $2.977\times10^{-5}$ & $-0.78\,\%$ & 5 \\
\bottomrule
\end{tabular}
\end{table}

A practical consequence for the hardware is that the blocks need not be
contiguous: each requires only $N+M$ consecutive bits, and independent
blocks give a slightly better estimate ($-0.30\,\%$ error with 20\,000
independent chunks). The capture buffer in the FPGA therefore holds a
single chunk of a few kilobits, filled at ring speed and read out slowly,
instead of megabits of contiguous samples.

\section{Transistor-level bench and its resolution}
\label{sec:spice}

The entropy source is the only block of the engine tied to the
technology, and in an integrated circuit it is designed at transistor
level. A SPICE bench was built around a five-stage ring with an enabling
NAND gate, with every technology-dependent parameter (device models,
widths, minimum length, supply) confined to a header so that moving from
the generic placeholder models used here to a foundry \ac{PDK} changes
one block of the netlist. Jitter is extracted from the threshold-crossing
instants of the output by the second difference
$D(M) = t_{k+2M} - 2t_{k+M} + t_k$, i.e.\ the Allan variance applied to
phase data, whose expectation $2M\sigma^2$ is free of the bias that a
simple difference over overlapping windows carries (measured at 5\,\%),
and which is the method used by Benea \emph{et al.} to characterize rings
both in Artix-7 and in 28~nm silicon~\cite{benea_2024_flicker}. Estimates
at different $M$ are combined with inverse-variance weights proportional
to the number of independent windows; without that weighting the
dispersion stayed at 15\,\% regardless of the simulated length.

The bench measures its own resolution first: with all noise sources
disabled, the residual jitter is the numerical noise of the simulator.
The first attempt gave 0.33~ps, of the same order as the jitter to be
measured; it was entirely the start-up transient of the oscillator, and
discarding the first twenty cycles brings the floor to \textbf{1.38~fs},
three orders of magnitude below the signal. Only then was noise injected.
Jitter scales with the square root of the injected spectral density
($\times2.01$ measured for $\times4$ in density, against $\times2.00$
expected) up to the point where the r.m.s.\ noise current reaches
about 80\,\% of the stage drive current; there the small-signal regime
breaks down, the mean period shifts, and the scaling law fails by a
factor of five. A validity criterion follows for every campaign,
including those on a real \ac{PDK}: keep injected noise below $\sim40\,\%$
of the stage current and verify that the mean period does not move when
noise is enabled. With the placeholder models the ring oscillates at
4.3~GHz, which is not physical; the numbers that enter the thesis come
from the foundry models, where transient noise analysis injects the
device noise itself.

\section{Design target and divider selection}
\label{sec:kd}

With $\sigma/T$ measured per period, $Q_{\mathrm{raw}} = (\sigma/T)^2$ and
the divider follows from $K_D = \lceil Q_{\mathrm{target}} /
Q_{\mathrm{raw}} \rceil$. For the range of jitter expected in a 28~nm
process the divider lies between $3\times10^4$ and $5\times10^5$ and the
throughput per ring between 1 and 20~kbit/s. The model reproduces the
published design point of Fischer and Lubicz ($K_D \approx 430\,000$ for
$\sigma = 5.01$~ps over $T = 7.81$~ns)~\cite{fischer_2014_embedded}. Such
rates are sufficient here: seeding the deterministic generator with
256~bits of full entropy requires about 512 raw bits at $H_\infty \approx
0.98$, i.e.\ between 0.03 and 0.5~s per seed, and the engine issues a seed
only at start-up and at reseed intervals. If more throughput were ever
needed, the clean route is several independent \acp{ERO}, each with its
own bound, rather than a multi-ring XOR whose independence cannot be
demonstrated~\cite{bochard_2010_true}.

\chapter{Digital Engine}
\label{ch:ch4_engine}

% [DRAFT] Source material: see docs/ in the repository.

\section{Architecture and the single-primitive decision}

% TODO

\section{AES-128 encryption core}

% TODO

\section{AES-CMAC}

% TODO

\section{CTR\_DRBG with derivation function}

% TODO

\section{Health tests}

% TODO

\section{Capture path and clock-domain crossings}

% TODO

\section{I2C register interface and command sequencer}

% TODO


\chapter{Verification}
\label{ch:ch5_verification}

% [DRAFT v1] Source material: rtl/tb/*, rtl/run_tb.py, analysis/ctr_drbg_ref.py,
% analysis/verifica_top.py, firmware/verificador/, docs/mapa_registros.md.
% FPGA bring-up section pending hardware.

\section{Strategy: vectors, simulation, independent verifier}
\label{sec:verif_strategy}

A cryptographic engine that is verified only against itself proves
nothing: an error in the S-box would be reproduced identically in the
model used to check it. The verification of this work therefore rests on
three legs that do not share an implementation.

\begin{enumerate}
\item \textbf{Official vectors.} Every primitive is checked against the
test vectors published with its standard: FIPS~197 and SP~800-38A for
AES, RFC~4493 and the NIST example file for CMAC, and the CAVP response
files for CTR\_DRBG (with and without derivation function, with and
without reseed). These are known-answer tests and they settle
correctness of the algorithm.
\item \textbf{Independent reference models.} A bit-exact model of the
generator was written in Python with its own AES, derived from the
standard without looking at the VHDL. It is validated against the CAVP
vectors first, and only then used to generate expected values for the
fixed input layout of the hardware and to recompute the tags of the
system-level bench. Two implementations that agree are worth more than
one that passes.
\item \textbf{An independent verifier in the loop.} On hardware the
ESP32 recomputes with mbedtls, and with its own AES accelerator, every
result the engine produces, so that the final check is performed by a
third, unrelated implementation over a real bus.
\end{enumerate}

Two working rules apply throughout, both inherited from the author's
previous experimental work: the output of a compilation, simulation or
measurement is never filtered before being read in full --- a filter
hides exactly the warning that mattered --- and failures are recorded
whole, never discarded. Twice in this work a filter on the simulator's
log hid a result, and both times the fix was to remove the filter.

\section{Reference models in Python}
\label{sec:refmodels}

\texttt{ctr\_drbg\_ref.py} implements AES-128 encryption, the block-cipher
derivation function, CBC-MAC and CTR\_DRBG with and without derivation
function, in pure Python and from the text of SP~800-90A. The AES asserts
the two FIPS~197 vectors at import; the DRBG is exercised against the CAVP
response files following the CAVS procedure (instantiate, optional reseed,
two generate calls, compare the second): \textbf{960 AES-128 cases, all
passing}. One defect was found in the file parser rather than in the
model --- a section header must close the pending test case before the
parameters change, otherwise the last case of every parameter block is
evaluated with the next block's parameters --- and it is worth recording
because its symptom, a single failure at count~14 of each block, looked
exactly like an algorithmic error.

A second script, \texttt{gen\_drbg\_vectores.py}, drives the validated
model through the sequence the hardware supports (instantiate with 32
bytes of entropy and 16 of nonce, generate twice, reseed, generate) and
emits a VHDL package of expected values; if the input layout of the
hardware ever changes, that script is the only place to change it. A
third, \texttt{verifica\_top.py}, recomputes AES-CMAC from the same AES and
validates itself against RFC~4493 before checking anything.

\section{Testbenches and results}
\label{sec:testbenches}

All VHDL is simulated with GHDL~6 under a single runner that analyses the
sources in dependency order and executes every bench, printing the
simulator's output unfiltered. Table~\ref{tab:tb} summarizes the seven
benches and their state at the time of writing.

\begin{table}[htbp]
\centering
\caption{Testbenches, what they check and their result (GHDL 6.0, VHDL-2008).}
\label{tab:tb}
\begin{tabular}{@{}llc@{}}
\toprule
Bench & Checks & Result \\
\midrule
Ring oscillator & start/stop, $f = 1/(2N t_{pd})$, zero floor, $\sigma_T = \sqrt{2N}\,\sigma_{\mathrm{stage}}$ & 4/4 \\
I2C slave & addressing, foreign address, multi-byte write, repeated-start read & 14/14 \\
AES-128 core & FIPS~197 B, C.1; SP~800-38A F.1.1 (3 blocks); 1000-block chain & 6/6 \\
AES-CMAC & RFC~4493 examples 1--4; repeatability & 5/5 \\
CTR\_DRBG & instantiate, 2$\times$generate, reseed, counter, re-instantiate, refuse unseeded & 13/13 \\
Health tests & healthy source, stuck, 70\,\% biased, stuck-at-complement, sticky/clear & 10/10 \\
Whole engine & ESP32 sequence over I2C: id, seed, key, 2 challenges, 2nd key, capture, key hidden & 16/16 \\
\bottomrule
\end{tabular}
\end{table}

Some of what the benches found deserves to be stated, because a
verification chapter that reports only passes has not verified much.

\paragraph{Ring model.} The LUT primitives used to simulate the ring
were first modeled with a strict unknown-propagation rule: any unknown
input gave an unknown output. At the start of a simulation every signal
is unknown, so the ring never left that state and never started. A real
NAND with its enable input at zero outputs one regardless of the other
input; the model now evaluates every combination compatible with the
known inputs and resolves the output when they all agree. Without this
the entire system-level bench would have been impossible.

\paragraph{I2C slave.} Two genuine design errors: the register pointer
did not advance on multi-byte writes, so consecutive bytes overwrote the
same address; and on reads the most significant bit was driven one clock
edge too late, shifting every byte. The first was fixed by incrementing
the pointer one state after the write strobe, so the address is still
valid during the strobe; the second by loading the shift register and
driving its first bit on the same falling edge that releases the
acknowledge. A third defect was in the bench itself: initializing the
register memory from the stimulus process created a second driver on the
same signal and every conflicting bit resolved to unknown.

\paragraph{Health tests.} The ``healthy'' source of the first bench, a
shift register seeded with a single one, started with a run of 27 zeros
and a first window at 76\,\%; the detector fired. The detector was right.

\paragraph{Key leak through the random register.} The most serious
finding did not come from a bench but from reading the design again
before handing it over. The command that generated the key wrote the
full 256-bit output of the DRBG into the random-output register, which is
readable without the test pin; the first 128 of those bits \emph{were} the
key. The system bench even asserted that equality as a ``check''. The fix
separates the two: key generation leaves nothing readable, and a new
command fills the random register from an independent generation, after
the DRBG has updated its state. The bench now asserts the opposite, that
the register contains the key in neither half, and the Python
cross-check does the same.

\paragraph{Whole engine.} With the rings at their physical speed the
system bench, which drives almost a millisecond of I2C traffic, did not
finish in twenty minutes of CPU: every ring transition costs two random
draws in the jitter model. In that bench alone the stage delay is set ten
times larger, with the same relative jitter, and this is stated in the
bench. The sixteen checks pass, and the key, the two challenges and their
tags are written to a file that \texttt{verifica\_top.py} checks: both
tags equal the CMAC recomputed by the independent Python AES, the
readable random-output register contains the key in neither of its
halves, and two generations give different keys. This is the first point at which a key born from
ring-oscillator noise, conditioned, generated and used to authenticate,
is confirmed by an implementation that shares nothing with the VHDL.

\section{ESP32 verifier firmware}
\label{sec:firmware}

The firmware is an ESP-IDF~5 application for the classic ESP32 that
talks to the engine through the register map of Appendix~\ref{app:regs}
and offers a console over the UART. Its commands map onto the
verification plan: \texttt{id} reads identification and status;
\texttt{sembrar} and \texttt{generar} issue the long commands and poll
the busy flag; in test mode \texttt{generar} also reads back the key
(provisioning) and checks that it equals the first sixteen bytes of the
generated block; \texttt{reto N} performs $N$ challenge-response rounds
with nonces from \texttt{esp\_random()}, recomputing each tag with
\texttt{mbedtls\_cipher\_cmac} and counting matches and mismatches, the
mismatches being printed in full; \texttt{crudo N} dumps $N$ chunks of
raw bits and \texttt{esp N} dumps blocks of the ESP32's own generator, both
in the framed, CRC32-protected format used in the author's previous work
so that the same statistical pipeline processes the two sources. CMAC is
not enabled by default in the ESP-IDF build of mbedtls and is switched on
in the project defaults; the console runs at 115\,200~baud, the only rate
that proved reliable with the board's USB bridge in earlier campaigns,
and the log level is set to errors only so that log lines cannot corrupt
the frames. The firmware builds with ESP-IDF~5.3.2 into a 255~kB image;
running it awaits the hardware.

\section{FPGA bring-up}
\label{sec:bringup}

% [PENDING HARDWARE] Vivado project, constraints (DONT_TOUCH, pblocks,
% ALLOW_COMBINATORIAL_LOOPS, set_disable_timing), first bitstream: do the
% rings oscillate and at what frequency; then the ESP32 in the loop.
Pending hardware. The constraints the rings need are known from the
literature and from the vendor documentation --- \texttt{DONT\_TOUCH} on
the cells and nets of each ring (unlike \texttt{KEEP}, it survives
implementation), \texttt{ALLOW\_COMBINATORIAL\_LOOPS} on one net of each
loop, \texttt{set\_disable\_timing} to break the loop for static timing
analysis, and one placement block per ring so that identical rings are
neither merged nor placed side by side --- and the first bitstream has a
single purpose: to read the frequency counters and confirm that every
ring oscillates. Nothing about entropy is claimed before that.

\chapter{Path to an Integrated Circuit}
\label{ch:ch6_asic}

% [DRAFT v1] Source material: docs/sintesis_resultados.md, docs/via_asic.md,
% synth/sintetiza.py, results/sintesis_sg13g2.csv.

The engine is designed to end up as a dedicated integrated circuit that
serves keys to a microcontroller, and the FPGA is its verification
vehicle, not its destination. This chapter takes the same VHDL that
passes the testbenches of Chapter~\ref{ch:ch5_verification}, synthesizes
it on an open standard-cell library, and states what a fabrication run
would cost and what would have to change in the entropy source. A
tape-out is explicitly out of scope of this thesis --- not for cost or
area, as will be seen, but for schedule --- and everything here is
labelled as what it is: a lower bound on area from logic synthesis, with
the growth factors to core area taken from published data.

\section{Standard-cell synthesis flow}
\label{sec:synthflow}

The flow is deliberately reproducible with free tools. Yosys reads the
VHDL through the GHDL front end, synthesizes each block flattened, maps
flip-flops and logic onto the liberty file of the target library and
reports cell count and area. The library is the open \ac{PDK} of IHP,
SG13G2, a 130~nm BiCMOS process with a maintained standard-cell library,
compiled SRAM macros, I/O cells and OpenROAD support, chosen because it is
the one open \ac{PDK} that is alive, European and reachable through the
same Europractice membership the university already holds. The gate
equivalent is read from the liberty file itself rather than from memory:
$1~\mathrm{GE} = $ area of \texttt{sg13g2\_nand2\_1} $= 7.2576~\mu\mathrm{m}^2$.
No clock constraint is applied: Yosys does not use one for area, so the
figures do not depend on a target frequency. The ring oscillators are not
synthesized; they are the one block designed at transistor level
(Chapter~\ref{ch:ch3_entropy}).

\section{Area per block in gate equivalents}
\label{sec:area}

Table~\ref{tab:synth} gives the result for each block. The digital
engine --- everything except the capture buffer --- is about
\textbf{62~kGE}, and AES is 51\,\% of it: the authenticator and the
generator each instantiate their own encryption core, 15.8~kGE apiece.
The capture buffer is reported for completeness but must not be read as
part of the product: its 46.7~kGE are a 4096-bit memory mapped onto
flip-flops because the flow has no SRAM macros. In the FPGA it occupies a
block RAM at no logic cost; in silicon it would be a compiled SRAM or,
more sensibly, a much smaller buffer, because it is a test path that is
disabled at production.

\begin{table}[htbp]
\centering
\caption{Logic synthesis on IHP SG13G2, flattened, no routing.
1~GE $=7.2576~\mu\mathrm{m}^2$. Lower bound on area \simu{}.}
\label{tab:synth}
\begin{tabular}{@{}lrrrr@{}}
\toprule
Block & Cells & Flip-flops & Area ($\mu\mathrm{m}^2$) & kGE \\
\midrule
AES-128 encryption core       & 10\,439 & 262   & 114\,771 & 15.81 \\
AES-CMAC (with its AES)        & 13\,210 & 909   & 169\,766 & 23.39 \\
CTR\_DRBG with df (with its AES) & 17\,953 & 2\,355 & 267\,262 & 36.83 \\
Health tests (RCT, APT)        & 568     & 85    & 8\,986   & 1.24 \\
ERO sampler and divider        & 117     & 24    & 2\,078   & 0.29 \\
Bit clock-domain crossing      & 11      & 7     & 408      & 0.06 \\
I2C slave                      & 310     & 63    & 5\,507   & 0.76 \\
\midrule
Capture buffer (4096-bit RAM as flip-flops) & 12\,250 & 4\,140 & 339\,057 & 46.72 \\
\bottomrule
\end{tabular}
\end{table}

The numbers are internally consistent, and that consistency is how an
error in the flow was caught: before flattening, the report returned for
the authenticator and for the generator the area of their AES
sub-module alone. An authenticator weighing exactly as much as a bare AES
is impossible, so the report, not the design, was wrong. After
flattening, the authenticator adds to the AES 647 flip-flops (sub-keys,
chaining state, key) and 7.6~kGE, and the generator adds 2093 flip-flops
and 21~kGE, about half of which are its registers.

\section{Optimizations and their expected savings}
\label{sec:opt}

The synthesis is naive on purpose --- the \ac{RTL} was written to be
verifiable, not small --- and it points at three optimizations, in order
of return on effort:

\begin{enumerate}
\item \textbf{One shared AES core.} The top-level sequencer already
serializes commands: the generator and the authenticator never run at the
same time. Sharing the core saves $\approx$15.8~kGE, a quarter of the
engine, for a trivial arbiter.
\item \textbf{Compact S-box datapath.} Twenty independent S-box tables
synthesize to 15.8~kGE for the whole core; area-optimized designs in the
literature range from 2.4~kGE for an 8-bit datapath at 226 cycles per
block~\cite{moradi2011pushing} to about 5~kGE for 32-bit datapaths. Speed
is irrelevant here, so a 32-bit datapath with four shared S-boxes is the
sensible point; estimated saving, 10~kGE per core.
\item \textbf{Fewer registers in the generator.} Of its 2355 flip-flops,
the captured entropy and nonce (384), the seed material (256), the output
(256) and the temporaries of the derivation function can be reduced by
consuming the seed as a stream and sharing temporaries; roughly 800
flip-flops, about 4~kGE.
\end{enumerate}

With the three, the digital engine falls to \textbf{$\approx$20~kGE},
which is the order of magnitude the state of the art anticipated (13 to
25~kGE for engines of this kind) and which, at the published growth
factor of $\approx$2.35 from cell area to core area for an AES on this same
library, is about 0.34~mm$^2$ of core. Each of these savings is stated as
an estimate; measuring them by synthesizing the optimized variants is
listed as future work, not claimed.

\section{What the ring oscillator needs in silicon}
\label{sec:ro_silicon}

The digital blocks port with the standard flow. The entropy source does
not, and the literature that has done the crossing --- notably a 2024
study that characterizes the same ring model on an Artix-7 and on a 28~nm
FD-SOI circuit~\cite{benea_2024_flicker} --- makes three points that the
chip design must absorb.

First, the ring is rebuilt from inverter cells rather than look-up tables,
and its jitter changes in value and in character: flicker noise gains
weight in newer nodes, and the study finds, against expectation, that
counting it as a legitimate source can raise throughput for the same
entropy. The stochastic model and the estimator of
Chapter~\ref{ch:ch3_entropy} carry over unchanged; the measured $\sigma$
does not. Second, process variation makes every die different, so the
frequency and jitter of each ring must be calibrated per chip; the
published all-digital designs that survive voltage and temperature
sweeps do it with a configurable ring and a tuning loop~\cite{yang_2016_jssc}.
Third, and most important for a security device, injecting a frequency on
the supply locks the rings and removes the jitter: a 500~mV signal on
the rail reduced the key space of a commercial secure microcontroller
from $2^{64}$ to about 3300~\cite{markettos_2009_frequency}, and a few
hundred microwatts of radiated field destroyed the statistics of a
fifty-ring generator~\cite{bayon_2012_contactless}. The same calibration
loop, together with supply filtering, is the published defence, and the
health tests of Section~\ref{sec:health} are the detector of last resort.
None of this is in the FPGA prototype; it is the work that a silicon
implementation adds.

\section{Cost and schedule of an academic shuttle}
\label{sec:shuttle}

Fabrication is more accessible than the size of the word suggests. The
University of Sevilla and IMSE-CNM are members of Europractice, whose
academic price list is public: a mini@sic block in UMC~180~nm costs
3\,430~\texteuro{}, one in TSMC~65~nm 3\,691~\texteuro{}, and IHP SG13G2 is
offered at 5\,110~\texteuro{}/mm$^2$ with a free multi-project run for
open-source designs under 2~mm$^2$~\cite{europractice_prices2026}. An engine of
0.3--0.4~mm$^2$ fits many times over in the smallest block of any of
them; the chip would be limited by its pads, not its logic. The
institute's own hardware-security group has already taken ring-oscillator
primitives from Xilinx~7-series FPGAs to characterized 65~nm silicon,
which makes the continuation of this work a matter of joining an
existing line rather than opening one.

What does not fit in a master's thesis is the schedule: from tape-out to
packaged samples takes months to over a year, the price does not include
packaging, and the characterization of the returned dies is a project in
itself. For comparison, a commercial secure element such as the
ATECC608B costs 0.90~\$ in single units and ships with a NIST entropy
source validation~\cite{microchip_atecc608b_ds}; two of its direct competitors
do not mention SP~800-90 anywhere in their datasheets. The honest
position of this work is therefore not to compete on price but on
\emph{transparency of the entropy}: every bit of a key produced by this
engine can be traced to a measured jitter and a stated bound, which is
precisely what a black-box part cannot offer.

\chapter{Results and Comparison}
\label{ch:ch7_results}

This chapter collects what the work establishes and separates it, without
softening, from what it does not. The engine has been modelled, written,
verified in simulation against official vectors and synthesized on a real
standard-cell library; it has not yet run on the \ac{FPGA}. Every number
below therefore carries its provenance tag, and the entropy campaign ---
the one result that requires silicon on the bench --- is stated as a
protocol to execute, not as data.

\section{Status against the objectives}
\label{sec:status}

Table~\ref{tab:status} maps the objectives of
Section~\ref{sec:objectives} onto evidence.

\begin{table}[htbp]
\centering
\caption{Objectives and the evidence supporting each.}
\label{tab:status}
\small
\begin{tabular}{@{}llp{7.4cm}@{}}
\toprule
Obj. & State & Evidence \\
\midrule
O1 & met    & Ring described at transistor level; noise swept in SPICE
              (Section~\ref{sec:spice_results}) \\
O2 & met    & Jitter extracted from the waveforms; estimator validated on
              synthetic data with known $Q$ \\
O3 & met    & Stochastic model closed: exact wrapped-Gaussian computation,
              design point $Q \ge 0.2286$ (Section~\ref{sec:opoint}) \\
O4 & met    & Engine written in portable VHDL-2008; seven testbenches pass
              against FIPS-197, SP 800-38B, RFC 4493 and \ac{CAVP} vectors \\
O5 & partial& Verifier firmware compiles and the protocol is closed; the
              board-level run is pending \\
O6 & pending& Requires the \ac{FPGA} on the bench
              (Section~\ref{sec:protocol}) \\
O7 & met    & Synthesis on IHP SG13G2, per-block area and optimization path
              (Chapter~\ref{ch:ch6_asic}) \\
O8 & partial& Structural comparison closed below; the numerical comparison
              against the ESP32 needs O6 \\
\bottomrule
\end{tabular}
\end{table}

The minimum defensible result declared in Chapter~\ref{ch:ch1_intro} ---
O1--O3 together with O5 --- is reached on its modelling and design side.
What O5 still lacks is the bring-up, not the design.

\section{Transistor-level noise of the ring}
\label{sec:spice_results}

The ring was simulated with a current noise source of controlled spectral
density injected at the switching node, sweeping the density and measuring
the standard deviation of the rising-edge crossings over 500 periods.
Table~\ref{tab:spice} gives the result for a five-stage ring with a
10~fF load.

\begin{table}[htbp]
\centering
\caption{Period jitter of the five-stage ring against injected noise
density \simu{}. $\sigma_{\mathrm{rel}} = \sigma/T$.}
\label{tab:spice}
\small
\begin{tabular}{@{}lrrrr@{}}
\toprule
Noise density (A$^2$/Hz) & Crossings & $T$ (ps) & $\sigma$ (fs) &
$\sigma_{\mathrm{rel}}$ \\
\midrule
none (numerical floor) & 500 & 230.30 & 1.38   & $6.0\times10^{-6}$ \\
$1\times10^{-20}$      & 501 & 230.37 & 1200   & $5.2\times10^{-3}$ \\
$4\times10^{-20}$      & 507 & 229.98 & 2416   & $1.1\times10^{-2}$ \\
$1.6\times10^{-19}$    & 509 & 227.77 & 25588  & $1.1\times10^{-1}$ \\
\bottomrule
\end{tabular}
\end{table}

Three readings, and the third is the one that matters.

\paragraph{The numerical floor is far enough below the signal.} With the
noise source switched off the measurement returns 1.38~fs, which is the
resolution of the simulation itself, not physics. It sits about 870 times
below the smallest jitter to be measured, so the floor does not
contaminate the sweep. Reaching it required care: the first attempt
returned 0.33~ps, of the same order as the quantity under study, and the
cause was the ring's start-up transient being included in the statistics.
Discarding the first twenty cycles removed it. A measurement floor that
coincides with the signal is indistinguishable from a result, which is
precisely why the floor was characterized before the sweep and not after.

\paragraph{The scaling law holds where it should.} Thermal jitter
accumulates as the square root of the noise density, so quadrupling the
density must double $\sigma$. From $1\times10^{-20}$ to
$4\times10^{-20}$ the measurement gives a factor of 2.01 against a
predicted 2.00. That agreement is the evidence that what is being measured
is the injected noise and not an artefact of the solver.

\paragraph{And it breaks at the top of the sweep, visibly.} At
$1.6\times10^{-19}$, sixteen times the reference density, $\sigma$ grows by
a factor of 21.3 where the law predicts 4.0, and the mean period drops from
230.4 to 227.8~ps. A zero-mean noise source cannot shift a mean period.
The point is therefore outside the regime the model describes --- the
perturbation is large enough to change the operating point of the
inverters --- and it is reported rather than dropped, because the boundary
of validity of a model is itself a result. Only the two central points are
used for anything downstream.

\section{Operating point: divider and throughput}
\label{sec:opoint}

The stochastic model fixes the target. For the PTG.2 thresholds of AIS
20/31 v3.0, the min-entropy criterion ($H_\infty \ge 0.98$) demands
$Q \ge 0.2286$, and is more demanding than the Shannon criterion
($H \ge 0.9998$, $Q \ge 0.2022$); the design therefore targets the former,
with $Q = 0.30$ as the working set-point, which yields
$H_\infty = 0.9951$. Two properties of the model were established rather
than assumed: the truncated series of equation~14
in~\cite{baudet_2011_security} \emph{overestimates} entropy and is thus
not a lower bound, so the exact wrapped-Gaussian computation is used
instead; and the first-harmonic derivation of min-entropy agrees with the
exact one to six digits over the whole range of interest.

Since $Q_{\mathrm{gen}} = K_D \cdot Q_{\mathrm{raw}}$, the divider follows
from the jitter that the campaign measures. Table~\ref{tab:opoint} gives
the operating points for the jitter range expected of this
technology~\cite{petura_2016_survey}, assuming RO2 at 500~MHz.

\begin{table}[htbp]
\centering
\caption{Divider and resulting throughput per ring for
$Q \ge 0.2286$ \est{}.}
\label{tab:opoint}
\small
\begin{tabular}{@{}rrrr@{}}
\toprule
$\sigma/T$ per period & $Q_{\mathrm{raw}}$ & $K_D$ & Throughput \\
\midrule
$7.1\times10^{-4}$ & $5\times10^{-7}$ & 457\,200 & 1.09~kbit/s \\
$1.4\times10^{-3}$ & $2\times10^{-6}$ & 114\,300 & 4.37~kbit/s \\
$2.8\times10^{-3}$ & $8\times10^{-6}$ & 28\,575  & 17.5~kbit/s \\
\bottomrule
\end{tabular}
\end{table}

These bracket the published design point of Fischer and
Lubicz~\cite{fischer_2014_embedded} ($K_D \approx 430\,000$ for
$\sigma = 5.01$~ps on $T = 7.81$~ns), which is the sanity check that the
model is being applied as its authors intended.

The architectural consequence is stated plainly: a single \ac{ERO}
delivers kbit/s, not Mbit/s. That is the price of demonstrable entropy,
and it is affordable here because the engine's product is a 256-bit seed
for the \ac{DRBG}, not a bit stream. More throughput means more
independent rings, each with its own bound, not a faster ring.

\section{Functional verification, consolidated}
\label{sec:verif_results}

Chapter~\ref{ch:ch5_verification} describes the strategy; this section
records the totals. Every figure in Table~\ref{tab:verif} was obtained by
running the suite, not by reading a previous log.

\begin{table}[htbp]
\centering
\caption{Verification totals \simu{}.}
\label{tab:verif}
\small
\begin{tabular}{@{}lll@{}}
\toprule
Suite & What it checks & Result \\
\midrule
VHDL testbenches (GHDL) & rings, \ac{ERO}, capture, health tests,
  \ac{I2C}, \ac{AES}, \ac{CMAC}, \ac{CTRDRBG}, top level & 7/7 pass \\
\texttt{ctr\_drbg\_ref.py} & \ac{CAVP} vectors, with and without
  derivation function & 960/960 \\
\texttt{ero\_model.py} & stochastic model, exactness and bounds & 24/24 \\
\texttt{jitter\_estimator.py} & jitter recovery on synthetic data of
  known $Q$ & 20/20 \\
\texttt{cortes\_salud.py} & health-test cutoffs against SP 800-90B
  formulas & 6/6 \\
\texttt{verifica\_top.py} & cross-check of the top-level bench with an
  independent \ac{AES} & 5/5 \\
\bottomrule
\end{tabular}
\end{table}

Two of these deserve a sentence. The 960 \ac{CAVP} cases are the
strongest single piece of evidence in the work: the generator is not
merely self-consistent, it reproduces the outputs that \ac{NIST} publishes
for \ac{CTRDRBG} with AES-128, in both the derivation-function and the
raw-entropy configurations. And the health-test cutoffs, $C_{RCT} = 22$
and $C_{APT} = 596$, had been computed by hand from the standard and were
never checked; deriving them again from the formulas of SP
800-90B~\cite{turan_2018_sp80090b} confirms both. The same check
quantifies the cost of the complement test added to the \ac{APT}, which is
an addition of this work and not of the standard: it raises the
false-alarm rate per window from $2^{-20.15}$ to $2^{-20.14}$, a 0.8\,\%
increase, because the complementary tail lies 5.7 standard deviations
away. The cutoff therefore remains valid as it stands and the addition is
free in practice.

\section{Comparison with commercial secure elements}
\label{sec:se_comparison}

The engine occupies the same functional slot as a commercial secure
element: a small companion device that generates and holds keys for a host
microcontroller over a two-wire bus. Table~\ref{tab:se} places it beside
the three parts that dominate the \ac{I2C} secure-element market.

\begin{table}[htbp]
\centering
\caption{The engine against commercial secure elements. Certification
data from the manufacturers' datasheets.}
\label{tab:se}
\small
\begin{tabular}{@{}lllll@{}}
\toprule
 & ATECC608B & SE050 & OPTIGA Trust M & This work \\
 & \cite{microchip_atecc608b_ds} & \cite{nxp_se050_ds} &
   \cite{infineon_optiga_trustm_ds} & \\
\midrule
Bus            & \ac{I2C} & \ac{I2C} & \ac{I2C} & \ac{I2C} \\
Certification  & FIPS 140-3 L2 & CC EAL6+ & CC EAL6+ & none \\
Entropy source & undisclosed & undisclosed & undisclosed & \ac{ERO},
                 published \\
Entropy bound  & not published & not published & not published &
                 derived from measured jitter \\
Key export     & not possible & not possible & not possible &
                 test mode only \\
Inspectable    & no & no & no & yes, \ac{RTL} released \\
\bottomrule
\end{tabular}
\end{table}

The honest reading of that table has two halves.

On capability and assurance the commercial parts win, and it is not close.
Two of them carry CC EAL6+ evaluations, which represent an amount of
third-party scrutiny --- including physical attack resistance --- that no
academic project reproduces. They offer asymmetric cryptography, certified
provisioning flows and manufacturing support. This work offers symmetric
key generation and authentication, and carries no certification at all.

On what the thesis is actually about, the comparison inverts. In all three
commercial devices the entropy source is undisclosed: the datasheets state
that a \ac{TRNG} is present, not what it is, and no min-entropy bound
derived from a measured physical parameter is published for any of them.
Their randomness is trusted because the vendor and the evaluation scheme
are trusted. The engine described here makes the opposite trade: it has no
certificate, and it can show where each bit of entropy comes from, what
physical measurement bounds it, and what the bound becomes when the
temperature or the divider changes. That is the same distinction that the
author's earlier characterization of the ESP32 generator ran into from the
other side --- a black box can be tested statistically but cannot be
bounded --- and it is the reason this engine was built as a white box.

Which of the two is preferable depends on the question being asked. For
shipping a product, the certified part. For knowing what the randomness in
that product is worth, only the open one answers.

\section{Uncertainty budget}
\label{sec:uncertainty}

Four quantities carry the conclusions, and each has a different and
declared kind of uncertainty.

\begin{description}
\item[Simulated jitter.] Resolution floor 1.38~fs \simu{}, about 870 times
below the smallest quantity measured. The scaling check against the
expected square-root law agrees to 0.5\,\% in the valid regime, and the
regime boundary is documented in Section~\ref{sec:spice_results}. The
figure depends on the noise density assumed for the technology, which is
itself \est{}: what the simulation establishes is the \emph{method} and
the ring's response, not the absolute jitter of a fabricated device.
\item[Recovered $Q$.] The estimator reproduces a known $Q$ with a bias
below 2\,\% and a dispersion of 4.0\,\% over eight seeds with 20\,000
blocks \simu{}. Without the corrections applied here a plain straight-line
fit biases the result by $-18$\,\%, so the correction is not cosmetic.
\item[Area.] A lower bound \simu{}: cells only, flattened, without routing,
fill or clock tree. Published growth factors from cell area to core area
on this library are $\approx 2.35$ for an \ac{AES}, and the optimization
savings of Chapter~\ref{ch:ch6_asic} are \est{} and explicitly not claimed
as measured.
\item[Min-entropy of the built device.] Not yet bounded. The model, the
estimator and the cutoffs are ready; the number requires the campaign of
Section~\ref{sec:protocol}. No value is offered here, and none should be
inferred from the simulated jitter.
\end{description}

\section{The pending campaign}
\label{sec:protocol}

The one result that silicon must supply is the min-entropy of the rings as
built, and the procedure is fixed in advance so that it cannot be tuned
after seeing the data.

\begin{enumerate}
\item \textbf{Confirm the rings oscillate.} Read the frequency of each
ring with the embedded counter before trusting any entropy measurement. A
ring removed by the synthesizer, or two rings that have locked to each
other, would otherwise be indistinguishable from a working source.
\item \textbf{Measure the jitter} with $K_D = 1$ by the method of Fischer
and Lubicz~\cite{fischer_2014_embedded}, sweeping the pair distance $M$
and fitting the linear region, with at least three rings so that the
global component separates from the local one~\cite{lubicz_2024_jitter}.
Fit $\sigma^2(t) = at + bt^2$ and keep only the thermal term.
\item \textbf{Fix $K_D$} from the measured jitter for $Q \ge 0.2286$ with
margin, and recompute the health-test cutoffs for the measured $H_\infty$
rather than the assumed 0.98, using \texttt{cortes\_salud.py}.
\item \textbf{Collect raw bits} at that operating point and run the
\ac{NIST} SP 800-90B estimators, contrasting the empirical estimate with
the model's bound. Agreement supports the model; disagreement is a result
about the model's hypotheses, and either way it is reported.
\item \textbf{Repeat across placement} --- rings in separated regions
versus adjacent ones --- to quantify the coupling that the literature
warns about, which is the measurement the architecture of this engine was
chosen to make possible.
\end{enumerate}

Until step 4 closes, the engine's claim is precise and limited: the chain
from physical noise to key material is designed, implemented and verified
against the standards, and the bound it will produce is computable from a
measurement that has not yet been taken.

\chapter{Conclusions and Future Work}
\label{ch:ch8_concl}

% [DRAFT] Source material: see docs/ in the repository.

\section{Conclusions}

% TODO

\section{Contributions}

% TODO

\section{Future work}

% TODO



% =========================================================================
\bibliographystyle{IEEEtran}
\bibliography{references}

% =========================================================================
\appendix
\chapter{Register Map}
\label{app:regs}
% TODO from docs/mapa_registros.md

\chapter{Official Test Vectors Used}
\label{app:vectors}
% TODO

\chapter{Reproduction Guide}
\label{app:repro}
% TODO


\end{document}
